% Keys: what a document sets around a diagram, in and out of groups, and what units and fonts mean there.
\documentclass{article}
\usepackage{amsmath,amssymb}
\usepackage{tikz-cd}
\pdfpagewidth=5in \pdfpageheight=5in
%%%% case cdset-every-arrow-append
\tikzcdset{every arrow/.append style={red}}
\begin{tikzcd} A \arrow[r] & B \end{tikzcd}
%%%% case cdset-every-arrow-replace
\tikzcdset{every arrow/.style={red}}
\begin{tikzcd} A \arrow[r] & B \end{tikzcd}
%%%% case cdset-every-arrow-prefix
\tikzcdset{every arrow/.prefix style={red}}
\begin{tikzcd} A \arrow[r] & B \end{tikzcd}
%%%% case cdset-every-arrow-group
{\tikzcdset{every arrow/.append style={red}}
\begin{tikzcd} A \arrow[r] & B \end{tikzcd}}
\begin{tikzcd} A \arrow[r] & B \end{tikzcd}
%%%% case cdset-arrows
\tikzcdset{arrows={red}}
\begin{tikzcd} A \arrow[r] & B \end{tikzcd}
%%%% case cdset-arrows-twice
\tikzcdset{arrows={red}}\tikzcdset{arrows={thick}}
\begin{tikzcd} A \arrow[r] & B \end{tikzcd}
%%%% case cdset-arrows-direction
\tikzcdset{arrows={r}}
\begin{tikzcd} A \arrow[] & B \end{tikzcd}
%%%% case cdset-arrows-label
\tikzcdset{arrows={"f"}}
\begin{tikzcd} A \arrow[r] & B \end{tikzcd}
%%%% case cdset-arrows-unknown
\tikzcdset{arrows={nonsensekey}}
\begin{tikzcd} A \arrow[r] & B \end{tikzcd}
%%%% case cdset-arrows-tikzcd-key
\tikzcdset{arrows={hook}}
\begin{tikzcd} A \arrow[r] & B \end{tikzcd}
%%%% case cdset-labels
\tikzcdset{labels={red}}
\begin{tikzcd} A \arrow[r, "f"] & B \end{tikzcd}
%%%% case cdset-labels-font
\tikzcdset{labels={font=\large}}
\begin{tikzcd} A \arrow[r, "f"] & B \end{tikzcd}
%%%% case cdset-every-label-style
\tikzcdset{every label/.style={red}}
\begin{tikzcd} A \arrow[r, "f"] & B \end{tikzcd}
%%%% case cdset-every-label-append
\tikzcdset{every label/.append style={font=\normalsize}}
\begin{tikzcd} A \arrow[r, "f"] & B \end{tikzcd}
%%%% case cdset-every-label-inner-sep
\tikzcdset{every label/.append style={inner sep=5pt}}
\begin{tikzcd} A \arrow[r, "f"] & B \end{tikzcd}
%%%% case cdset-every-cell
\tikzcdset{every cell/.append style={draw}}
\begin{tikzcd} A \arrow[r] & B \end{tikzcd}
%%%% case cdset-every-cell-style
\tikzcdset{every cell/.style={draw}}
\begin{tikzcd} A \arrow[r] & B \end{tikzcd}
%%%% case cdset-every-matrix-append
\tikzcdset{every matrix/.append style={inner sep=5pt}}
\begin{tikzcd} A \arrow[r] & B \end{tikzcd}
%%%% case cdset-every-matrix-line-width
\tikzcdset{every matrix/.append style={line width=2pt}}
\begin{tikzcd} A \arrow[r] & B \end{tikzcd}
%%%% case cdset-every-matrix-draw
\tikzcdset{every matrix/.append style={draw}}
\begin{tikzcd} A \arrow[r] & B \end{tikzcd}
%%%% case cdset-every-matrix-font
\tikzcdset{every matrix/.append style={font=\large}}
\begin{tikzcd} A \arrow[r] & B \end{tikzcd}
%%%% case cdset-every-matrix-style
\tikzcdset{every matrix/.style={inner sep=5pt}}
\begin{tikzcd} A \arrow[r] & B \end{tikzcd}
%%%% case cdset-every-diagram-append
\tikzcdset{every diagram/.append style={row sep=2cm}}
\begin{tikzcd} A \arrow[d] \\ B \end{tikzcd}
%%%% case cdset-every-diagram-style
\tikzcdset{every diagram/.style={row sep=2cm}}
\begin{tikzcd} A \arrow[d] \\ B \end{tikzcd}
%%%% case cdset-every-diagram-baseline
\tikzcdset{every diagram/.append style={baseline=(current bounding box.south)}}
x\begin{tikzcd} A \arrow[d] \\ B \end{tikzcd}
%%%% case cdset-diagrams
\tikzcdset{diagrams={row sep=2cm}}
\begin{tikzcd} A \arrow[d] \\ B \end{tikzcd}
%%%% case cdset-cells
\tikzcdset{cells={nodes={draw}}}
\begin{tikzcd} A \arrow[r] & B \end{tikzcd}
%%%% case cdset-row-sep
\tikzcdset{row sep=2cm}
\begin{tikzcd} A \arrow[d] \\ B \end{tikzcd}
%%%% case cdset-row-sep-normal
\tikzcdset{row sep/normal=2cm}
\begin{tikzcd} A \arrow[d] \\ B \end{tikzcd}
%%%% case cdset-row-sep-normal-group
{\tikzcdset{row sep/normal=2cm}
\begin{tikzcd} A \arrow[d] \\ B \end{tikzcd}}
\begin{tikzcd} A \arrow[d] \\ B \end{tikzcd}
%%%% case cdset-column-sep-normal
\tikzcdset{column sep/normal=2cm}
\begin{tikzcd} A \arrow[r] & B \end{tikzcd}
%%%% case cdset-row-sep-normal-list
\tikzcdset{row sep/normal={between origins, 2cm}}
\begin{tikzcd} A \arrow[d] \\ B \end{tikzcd}
%%%% case cdset-row-sep-huge
\tikzcdset{row sep/huge=5cm}
\begin{tikzcd}[row sep=huge] A \arrow[d] \\ B \end{tikzcd}
%%%% case cdset-row-sep-macro
\def\q{2cm}\tikzcdset{row sep/normal=\q}
\begin{tikzcd} A \arrow[d] \\ B \end{tikzcd}
%%%% case cdset-row-sep-macro-changed
\def\q{2cm}\tikzcdset{row sep/normal=\q}
\def\q{4cm}
\begin{tikzcd} A \arrow[d] \\ B \end{tikzcd}
%%%% case cdset-row-sep-em
\tikzcdset{row sep/normal=3em}
\Large
\begin{tikzcd} A \arrow[d] \\ B \end{tikzcd}
%%%% case cdset-row-sep-new
\tikzcdset{row sep/mine/.initial=3cm}
\begin{tikzcd}[row sep=mine] A \arrow[d] \\ B \end{tikzcd}
%%%% case cdset-row-sep-undefined
\begin{tikzcd}[row sep=mine] A \arrow[d] \\ B \end{tikzcd}
%%%% case cdset-sep
\begin{tikzcd}[sep=large] A \arrow[d] \arrow[r] & B \\ C \end{tikzcd}
%%%% case cdset-column-row-scriptsize
\begin{tikzcd}[row sep=scriptsize, column sep=tiny] A \arrow[d] \arrow[r] & B \\ C \end{tikzcd}
%%%% case cdset-sep-tikz
\begin{tikzcd}[/tikz/row sep=1cm, /tikz/column sep=2cm] A \arrow[d] \arrow[r] & B \\ C \end{tikzcd}
%%%% case cdset-restored
{\tikzcdset{every arrow/.style={red}}}
\begin{tikzcd} A \arrow[r] & B \end{tikzcd}
%%%% case cdset-style-restored-same
\let\saved\relax
\makeatletter
\expandafter\let\expandafter\saved\csname pgfk@/tikz/commutative diagrams/every arrow/.@cmd\endcsname
\tikzcdset{every arrow/.style={red}}
\expandafter\let\csname pgfk@/tikz/commutative diagrams/every arrow/.@cmd\endcsname\saved
\makeatother
\begin{tikzcd} A \arrow[r] & B \end{tikzcd}
%%%% case cdset-style-redefined-same-text
\tikzcdset{every arrow/.style={/tikz/line width=rule_thickness,draw}}
\begin{tikzcd} A \arrow[r] & B \end{tikzcd}
%%%% case cdset-after-diagram
\begin{tikzcd} A \arrow[r] & B \end{tikzcd}
\tikzcdset{every arrow/.append style={red}}
\begin{tikzcd} A \arrow[r] & B \end{tikzcd}
\tikzcdset{every arrow/.append style={thick}}
\begin{tikzcd} A \arrow[r] & B \end{tikzcd}
%%%% case cdset-math-mode
\tikzcdset{math mode=false}
\begin{tikzcd} A \arrow[r, "f"] & $B$ \end{tikzcd}
%%%% case cdset-math-mode-option
\begin{tikzcd}[math mode=false] A \arrow[r, "f"] & $B$ \end{tikzcd}
%%%% case cdset-math-mode-group
\begin{tikzcd}[math mode=false] {$A$} \arrow[r, "$f$"] & text \end{tikzcd}
%%%% case cdset-math-mode-toggle
{\tikzcdset{math mode=false}}\tikzcdset{math mode=true}
\begin{tikzcd} A \arrow[r, "f"] & B \end{tikzcd}
%%%% case cdset-background-color
\tikzcdset{background color=yellow}
\begin{tikzcd} A \arrow[rd] & B \arrow[ld, crossing over] \\ C & D \end{tikzcd}
%%%% case cdset-background-macro
\def\q{yellow}\tikzcdset{background color=\q}
\begin{tikzcd} A \arrow[rd] & B \arrow[ld, crossing over] \\ C & D \end{tikzcd}
%%%% case cdset-clearance
\tikzcdset{crossing over clearance=3ex}
\begin{tikzcd} A \arrow[rd] & B \arrow[ld, crossing over] \\ C & D \end{tikzcd}
%%%% case cdset-cramped
\tikzcdset{cramped}
\begin{tikzcd} A \arrow[r] & B \end{tikzcd}
%%%% case cdset-cramped-option
\begin{tikzcd}[cramped] A \arrow[r] & B \end{tikzcd}
%%%% case cdset-new-style
\tikzcdset{q/.style={red}}
\begin{tikzcd}[q] A \arrow[r] & B \end{tikzcd}
%%%% case cdset-new-style-tikz
\tikzset{q/.style={red}}
\begin{tikzcd}[q] A \arrow[r] & B \end{tikzcd}
%%%% case cdset-new-style-text
\tikzset{q/.style={font=\large}}
\begin{tikzcd}[q] A \arrow[r] & B \end{tikzcd}
%%%% case tikzset-every-node
\tikzset{every node/.style={draw}}
\begin{tikzcd} A \arrow[r, "f"] & B \end{tikzcd}
%%%% case tikzset-every-node-group
{\tikzset{every node/.style={draw}}
\begin{tikzcd} A \arrow[r, "f"] & B \end{tikzcd}}
\begin{tikzcd} A \arrow[r, "f"] & B \end{tikzcd}
%%%% case tikzset-every-node-append
\tikzset{every node/.append style={font=\large}}
\begin{tikzcd} A \arrow[r, "f"] & B \end{tikzcd}
%%%% case tikzset-every-picture
\tikzset{every picture/.style={red}}
\begin{tikzcd} A \arrow[r, "f"] & B \end{tikzcd}
%%%% case tikzset-every-picture-append
\tikzset{every picture/.append style={baseline=1cm}}
\begin{tikzcd} A \arrow[r, "f"] & B \end{tikzcd}
%%%% case tikzset-every-path
\tikzset{every path/.style={red}}
\begin{tikzcd} A \arrow[r, "f"] & B \end{tikzcd}
%%%% case tikzset-every-path-prefix
\tikzset{every path/.prefix style={thick}}
\begin{tikzcd} A \arrow[r, "f"] & B \end{tikzcd}
%%%% case tikzset-every-to
\tikzset{every to/.style={red}}
\begin{tikzcd} A \arrow[r, "f"] & B \end{tikzcd}
%%%% case tikzset-every-matrix
\tikzset{every matrix/.style={draw}}
\begin{tikzcd} A \arrow[r, "f"] & B \end{tikzcd}
%%%% case tikzset-every-cell
\tikzset{every cell/.style={draw}}
\begin{tikzcd} A \arrow[r, "f"] & B \end{tikzcd}
%%%% case tikzset-every-rectangle-node
\tikzset{every rectangle node/.style={draw}}
\begin{tikzcd} A \arrow[r, "f"] & B \end{tikzcd}
%%%% case tikzset-every-text-node-part
\tikzset{every text node part/.style={red}}
\begin{tikzcd} A \arrow[r, "f"] & B \end{tikzcd}
%%%% case tikzset-every-label
\tikzset{every label/.style={red}}
\begin{tikzcd} A \arrow[r, "f"] & B \end{tikzcd}
%%%% case tikzset-every-edge
\tikzset{every edge/.style={red}}
\begin{tikzcd} A \arrow[r, "f"] & B \end{tikzcd}
%%%% case tikzset-every-scope
\tikzset{every scope/.style={red}}
\begin{tikzcd} A \arrow[r, "f"] & B \end{tikzcd}
%%%% case tikzset-every-coordinate-node
\tikzset{every coordinate node/.style={red}}
\begin{tikzcd} A \arrow[r, "f"] & B \end{tikzcd}
%%%% case tikzset-gt
\tikzset{>=stealth}
\begin{tikzcd} A \arrow[r] & B \end{tikzcd}
%%%% case tikzset-gt-used
\tikzset{>=stealth}
\begin{tikzcd} A \arrow[r, ->] \arrow[d, -to] & B \\ C \end{tikzcd}
%%%% case tikzset-line-width
\tikzset{line width=2pt}
\begin{tikzcd} A \arrow[r] & |[draw]| B \end{tikzcd}
%%%% case tikzset-thick
\tikzset{thick}
\begin{tikzcd} A \arrow[r] & |[draw]| B \end{tikzcd}
%%%% case tikzset-red
\tikzset{red}
\begin{tikzcd} A \arrow[r, "f"] & |[draw]| B \end{tikzcd}
%%%% case tikzset-text-red
\tikzset{text=red}
\begin{tikzcd} A \arrow[r, "f"] & B \end{tikzcd}
%%%% case tikzset-node-font
\tikzset{node font=\large}
\begin{tikzcd} A \arrow[r, "f"] & B \end{tikzcd}
%%%% case tikzset-font
\tikzset{font=\large}
\begin{tikzcd} A \arrow[r, "f"] & B \end{tikzcd}
%%%% case tikzset-inner-sep
\tikzset{inner sep=5pt}
\begin{tikzcd} A \arrow[r, "f"] & B \end{tikzcd}
%%%% case tikzset-outer-sep
\tikzset{outer sep=5pt}
\begin{tikzcd} A \arrow[r, "f"] & B \end{tikzcd}
%%%% case tikzset-minimum-size
\tikzset{minimum size=1cm}
\begin{tikzcd} A \arrow[r, "f"] & B \end{tikzcd}
%%%% case tikzset-text-height
\tikzset{text height=1cm}
\begin{tikzcd} A \arrow[r, "f"] & B \end{tikzcd}
%%%% case tikzset-text-depth
\tikzset{text depth=1cm}
\begin{tikzcd} A \arrow[r, "f"] & B \end{tikzcd}
%%%% case tikzset-text-width
\tikzset{text width=1cm}
\begin{tikzcd} A \arrow[r, "f"] & B \end{tikzcd}
%%%% case tikzset-align
\tikzset{align=center}
\begin{tikzcd} A \arrow[r, "f"] & B \end{tikzcd}
%%%% case tikzset-execute-begin-node
\tikzset{execute at begin node=Q}
\begin{tikzcd} A \arrow[r, "f"] & B \end{tikzcd}
%%%% case tikzset-execute-begin-picture
\tikzset{execute at begin picture=\typeout{x}}
\begin{tikzcd} A \arrow[r, "f"] & B \end{tikzcd}
%%%% case tikzset-execute-begin-cell
\tikzset{execute at begin cell=Q}
\begin{tikzcd} A \arrow[r, "f"] & B \end{tikzcd}
%%%% case tikzset-execute-end-cell
\tikzset{execute at end cell=Q}
\begin{tikzcd} A \arrow[r, "f"] & B \end{tikzcd}
%%%% case tikzset-execute-empty-cell
\tikzset{execute at empty cell=Q}
\begin{tikzcd} A \arrow[r, "f"] & \end{tikzcd}
%%%% case tikzset-shorten
\tikzset{shorten >=3pt}
\begin{tikzcd} A \arrow[r] & B \end{tikzcd}
%%%% case tikzset-shorten-start
\tikzset{shorten <=3pt}
\begin{tikzcd} A \arrow[r] & B \end{tikzcd}
%%%% case tikzset-rounded
\tikzset{rounded corners}
\begin{tikzcd} A \arrow[r] & B \end{tikzcd}
%%%% case tikzset-dashed
\tikzset{dashed}
\begin{tikzcd} A \arrow[r] & B \end{tikzcd}
%%%% case tikzset-opacity
\tikzset{opacity=.5}
\begin{tikzcd} A \arrow[r] & B \end{tikzcd}
%%%% case tikzset-scale
\tikzset{scale=2}
\begin{tikzcd} A \arrow[r] & B \end{tikzcd}
%%%% case tikzset-xscale
\tikzset{xscale=2}
\begin{tikzcd} A \arrow[r] & B \end{tikzcd}
%%%% case tikzset-rotate
\tikzset{rotate=30}
\begin{tikzcd} A \arrow[r] & B \end{tikzcd}
%%%% case tikzset-xshift
\tikzset{xshift=1cm}
\begin{tikzcd} A \arrow[r] & B \end{tikzcd}
%%%% case tikzset-x-unit
\tikzset{x=2cm,y=2cm}
\begin{tikzcd} A \arrow[r] & B \end{tikzcd}
%%%% case tikzset-baseline
\tikzset{baseline=1cm}
\begin{tikzcd} A \arrow[r] & B \end{tikzcd}
%%%% case tikzset-looseness
\tikzset{looseness=2}
\begin{tikzcd} A \arrow[r, bend left] & B \end{tikzcd}
%%%% case tikzset-bend-angle
\tikzset{bend angle=60}
\begin{tikzcd} A \arrow[r, bend left] & B \end{tikzcd}
%%%% case tikzset-auto
\tikzset{auto}
\begin{tikzcd} A \arrow[r, "f"] & B \end{tikzcd}
%%%% case tikzset-auto-right
\tikzset{auto=right}
\begin{tikzcd} A \arrow[r, "f"] & B \end{tikzcd}
%%%% case tikzset-swap
\tikzset{swap}
\begin{tikzcd} A \arrow[r, "f"] & B \end{tikzcd}
%%%% case tikzset-pos
\tikzset{pos=.2}
\begin{tikzcd} A \arrow[r, "f"] & B \end{tikzcd}
%%%% case tikzset-near-start
\tikzset{near start}
\begin{tikzcd} A \arrow[r, "f"] & B \end{tikzcd}
%%%% case tikzset-sloped
\tikzset{sloped}
\begin{tikzcd} A \arrow[r, "f"] & B \end{tikzcd}
%%%% case tikzset-draw
\tikzset{draw}
\begin{tikzcd} A \arrow[r, "f"] & B \end{tikzcd}
%%%% case tikzset-fill
\tikzset{fill=yellow}
\begin{tikzcd} A \arrow[r, "f"] & B \end{tikzcd}
%%%% case tikzset-nodes
\tikzset{nodes=draw}
\begin{tikzcd} A \arrow[r, "f"] & B \end{tikzcd}
%%%% case tikzset-nodes-red
\tikzset{nodes={red}}
\begin{tikzcd} A \arrow[r, "f"] & B \end{tikzcd}
%%%% case tikzset-column-1
\tikzset{column 1/.style={red}}
\begin{tikzcd} A \arrow[r, "f"] & B \end{tikzcd}
%%%% case tikzset-row-1
\tikzset{row 1/.style={red}}
\begin{tikzcd} A \arrow[r, "f"] & B \end{tikzcd}
%%%% case tikzset-row-1-column-2
\tikzset{row 1 column 2/.style={red}}
\begin{tikzcd} A \arrow[r, "f"] & B \end{tikzcd}
%%%% case tikzset-column-5-not-reached
\tikzset{column 5/.style={red}}
\begin{tikzcd} A \arrow[r, "f"] & B \end{tikzcd}
%%%% case tikzset-1-row-diagram
\tikzcdset{1-row diagram/.style={red}}
\begin{tikzcd} A \arrow[r, "f"] & B \end{tikzcd}
%%%% case tikzset-2-row-diagram
\tikzcdset{2-row diagram/.style={row sep=3cm}}
\begin{tikzcd} A \arrow[d, "f"] \\ B \end{tikzcd}
%%%% case tikzset-1-row-diagram-show
\tikzcdset{1-row diagram/.style={/tikz/baseline=1cm}}
\begin{tikzcd} A \arrow[r, "f"] & B \end{tikzcd}
%%%% case tikzset-matrix-anchor
\tikzset{matrix anchor=north}
\begin{tikzcd} A \arrow[r, "f"] & B \end{tikzcd}
%%%% case tikzset-anchor
\tikzset{anchor=north}
\begin{tikzcd} A \arrow[r, "f"] & B \end{tikzcd}
%%%% case tikzset-left
\tikzset{left}
\begin{tikzcd} A \arrow[r, "f"] & B \end{tikzcd}
%%%% case tikzset-to-path
\tikzset{to path={-- (\tikztotarget)}}
\begin{tikzcd} A \arrow[r, "f"] & B \end{tikzcd}
%%%% case tikzset-line-cap
\tikzset{line cap=round, line join=bevel, miter limit=5}
\begin{tikzcd} A \arrow[r, "f"] & B \end{tikzcd}
%%%% case tikzset-dash
\tikzset{dash pattern=on 2pt off 2pt, dash phase=1pt}
\begin{tikzcd} A \arrow[r, "f"] & B \end{tikzcd}
%%%% case tikzset-double
\tikzset{double}
\begin{tikzcd} A \arrow[r, "f"] & B \end{tikzcd}
%%%% case tikzset-double-distance
\tikzset{double distance=3pt}
\begin{tikzcd} A \arrow[r, Rightarrow] & B \end{tikzcd}
%%%% case tikzset-pattern
\usetikzlibrary{patterns}
\tikzset{pattern=dots}
\begin{tikzcd} A \arrow[r, "f"] & B \end{tikzcd}
%%%% case tikzset-transparency
\tikzset{transparency group}
\begin{tikzcd} A \arrow[r, "f"] & B \end{tikzcd}
%%%% case tikzset-handle-active-nodes
\tikzset{handle active characters in nodes}
\begin{tikzcd} A \arrow[r, "f"] & B \end{tikzcd}
%%%% case tikzset-handle-active-code
\tikzset{handle active characters in code}
\begin{tikzcd} A \arrow[r, "f"] & B \end{tikzcd}
%%%% case tikzset-handle-active-with-char
\tikzset{handle active characters in nodes}
\catcode`\;=13 \def;{\mathrm{z}}
\begin{tikzcd} A ; \arrow[r, "f"] & B \end{tikzcd}
\catcode`\;=12
%%%% case tikzset-handle-active-code-with-char
\tikzset{handle active characters in code}
\catcode`\;=13 \def;{\mathrm{z}}
\begin{tikzcd} A ; \arrow[r, "f"] & B \end{tikzcd}
\catcode`\;=12
%%%% case pgfkeys-direct-every-arrow
\pgfkeys{/tikz/commutative diagrams/every arrow/.append style={red}}
\begin{tikzcd} A \arrow[r] & B \end{tikzcd}
%%%% case pgfkeys-direct-every-node
\pgfkeys{/tikz/every node/.style={draw}}
\begin{tikzcd} A \arrow[r] & B \end{tikzcd}
%%%% case pgfkeys-direct-cmd
\expandafter\def\csname pgfk@/tikz/commutative diagrams/row sep/normal\endcsname{5cm}
\begin{tikzcd} A \arrow[d] \\ B \end{tikzcd}
%%%% case pgfkeys-direct-cd
\pgfkeys{/tikz/commutative diagrams/.cd, row sep/normal=3cm}
\begin{tikzcd} A \arrow[d] \\ B \end{tikzcd}
%%%% case pgfkeys-tikzcd-dir-in-options
\begin{tikzcd}[/tikz/commutative diagrams/.cd, row sep=3cm] A \arrow[d] \\ B \end{tikzcd}
%%%% case pgfkeys-tikz-dir-in-options
\begin{tikzcd}[/tikz/.cd, row sep=3cm] A \arrow[d] \\ B \end{tikzcd}
%%%% case pgfkeys-own-dir-in-options
\begin{tikzcd}[/my/.cd, row sep=3cm] A \arrow[d] \\ B \end{tikzcd}
%%%% case pgfkeys-tikzcd-dir-in-arrow
\begin{tikzcd} A \arrow[d, /tikz/commutative diagrams/.cd, hook] \\ B \end{tikzcd}
%%%% case pgfkeys-diagrams-in-options
\begin{tikzcd}[diagrams={row sep=3cm}] A \arrow[d] \\ B \end{tikzcd}
%%%% case pgfkeys-defaults
\pgfkeys{/tikz/commutative diagrams/row sep/.default=1cm}
\begin{tikzcd}[row sep] A \arrow[d] \\ B \end{tikzcd}
%%%% case pgfkeys-tikz-node-distance
\pgfkeys{/tikz/node distance=5cm}
\begin{tikzcd} A \arrow[d] \\ B \end{tikzcd}
%%%% case pgfkeys-pgf-stuff
\pgfkeys{/pgf/inner xsep=5pt, /pgf/inner ysep=5pt}
\begin{tikzcd} A \arrow[d] \\ B \end{tikzcd}
%%%% case pgfset-outer-sep
\pgfset{outer sep=5pt, minimum size=1cm}
\begin{tikzcd} A \arrow[d] \\ B \end{tikzcd}
%%%% case pgfset-fpu
\usepgflibrary{fpu}
\pgfkeys{/pgf/fpu}
\begin{tikzcd} A \arrow[d] \\ B \end{tikzcd}
\pgfkeys{/pgf/fpu=false}
%%%% case pgfset-number-format
\pgfkeys{/pgf/number format/fixed}
\begin{tikzcd} A \arrow[d] \\ B \end{tikzcd}
%%%% case pgfsetlinewidth-outside
\pgfsetlinewidth{3pt}
\begin{tikzcd} A \arrow[d] \\ |[draw]| B \end{tikzcd}
%%%% case pgfsetcolor-outside
\pgfsetcolor{red}
\begin{tikzcd} A \arrow[d] \\ B \end{tikzcd}
%%%% case pgfsetarrows-outside
\pgfsetarrows{stealth-stealth}
\begin{tikzcd} A \arrow[d] \\ B \end{tikzcd}
%%%% case pgfsetdash-outside
\pgfsetdash{{3pt}{1pt}}{0pt}
\begin{tikzcd} A \arrow[d] \\ B \end{tikzcd}
%%%% case pgftransform-outside
\pgftransformshift{\pgfpoint{1cm}{1cm}}\pgftransformscale{2}
\begin{tikzcd} A \arrow[d] \\ B \end{tikzcd}
%%%% case pgfmath-macros
\pgfmathsetmacro\q{3}
\begin{tikzcd}[row sep=\q cm] A \arrow[d] \\ B \end{tikzcd}
%%%% case pgfmath-register
\pgfmathsetlength\dimen0{1cm+1pt}
\begin{tikzcd}[row sep=\dimen0] A \arrow[d] \\ B \end{tikzcd}
%%%% case macro-pgfutil-redefined
\makeatletter
\let\pgfutil@ifnextchar\@ifnextchar
\makeatother
\begin{tikzcd} A \arrow[d] \\ B \end{tikzcd}
%%%% case macro-pgf-redefined
\renewcommand\pgfmatrixnextcell{Q}
\begin{tikzcd} A \arrow[r] & B \end{tikzcd}
%%%% case macro-tikz-redefined
\makeatletter
\def\tikz@do@fig{\relax}
\makeatother
\begin{tikzcd} A \arrow[r] & B \end{tikzcd}
%%%% case macro-tikzcd-internal-redefined
\makeatletter
\def\tikzcd@{\relax}
\makeatother
\begin{tikzcd} A \arrow[r] & B \end{tikzcd}
%%%% case macro-used-by-pgf
\renewcommand\ifthenelse{}
\renewcommand\normalbaselines{\baselineskip=30pt}
\begin{tikzcd} \parbox{1cm}{a b c d e f} \arrow[r] & B \end{tikzcd}
%%%% case macro-selectfont-redefined
\let\oldsf\selectfont \def\selectfont{\oldsf\typeout{sel}}
\begin{tikzcd} A \arrow[r] & B \end{tikzcd}
%%%% case macro-normalfont-redefined
\renewcommand\normalfont{\bfseries}
\begin{tikzcd} A \arrow[r] & B \end{tikzcd}
%%%% case macro-color-redefined
\renewcommand\color[1]{\typeout{col}}
\begin{tikzcd} A \arrow[r, red] & B \end{tikzcd}
%%%% case macro-leavevmode-redefined
\let\oldlv\leavevmode \def\leavevmode{Q\oldlv}
\begin{tikzcd} A \arrow[r] & B \end{tikzcd}
%%%% case macro-unskip-redefined
\let\oldun\unskip \def\unskip{Q\oldun}
\begin{tikzcd} A \arrow[r] & B \end{tikzcd}
%%%% case macro-ignorespaces-redefined
\let\oldig\ignorespaces \def\ignorespaces{Q\oldig}
\begin{tikzcd} A \arrow[r] & B \end{tikzcd}
%%%% case macro-relax-redefined
\let\oldrelax\relax \def\relax{Q}
\begin{tikzcd} A \arrow[r] & B \end{tikzcd}
%%%% case macro-begingroup-redefined
\let\oldbg\begingroup \def\begingroup{Q\oldbg}
\begin{tikzcd} A \arrow[r] & B \end{tikzcd}
%%%% case macro-endgroup-redefined
\let\oldeg\endgroup \def\endgroup{\oldeg Q}
\begin{tikzcd} A \arrow[r] & B \end{tikzcd}
%%%% case macro-aftergroup-redefined
\let\oldag\aftergroup \def\aftergroup{\oldag\relax\oldag}
\begin{tikzcd} A \arrow[r] & B \end{tikzcd}
%%%% case macro-expandafter-redefined
\let\oldex\expandafter \def\expandafter{\oldex}
\begin{tikzcd} A \arrow[r, "f"] & B \end{tikzcd}
%%%% case macro-the-redefined
\let\oldthe\the \def\the{\oldthe}
\begin{tikzcd} A \arrow[r, "f"] & B \end{tikzcd}
%%%% case macro-everymath-redefined
\newtoks\mytoks \let\everymath\mytoks
\begin{tikzcd} A \arrow[r, "f"] & B \end{tikzcd}
%%%% case macro-scriptstyle-redefined
\def\scriptstyle{\displaystyle}
\begin{tikzcd} A \arrow[r, "\sum_f"] & B \end{tikzcd}
%%%% case macro-scriptstyle-let-text
\let\scriptstyle\textstyle
\begin{tikzcd} A \arrow[r, "\sum_f"] & B \end{tikzcd}
%%%% case macro-hbox-redefined
\let\oldhb\hbox \def\hbox{Q\oldhb}
\begin{tikzcd} A \arrow[r] & B \end{tikzcd}
%%%% case macro-pdfliteral-redefined
\let\oldpl\pdfliteral \def\pdfliteral{\typeout{lit}\oldpl}
\begin{tikzcd} A \arrow[r] & B \end{tikzcd}
%%%% case macro-pdfcolorstack
\let\oldcs\pdfcolorstack \def\pdfcolorstack{\typeout{st}\oldcs}
\begin{tikzcd} A \arrow[r, red] & B \end{tikzcd}
%%%% case macro-begin-redefined-tracking
\let\oldbegin\begin \def\begin#1{\typeout{<#1>}\oldbegin{#1}}
\begin{tikzcd} A \arrow[r] & B \end{tikzcd}
%%%% case macro-end-redefined-tracking
\let\oldend\end \def\end#1{\typeout{<#1>}\oldend{#1}}
\begin{tikzcd} A \arrow[r] & B \end{tikzcd}
%%%% case macro-safe-redefined-mention
\renewcommand\bfseries{\arrow[r]}
\begin{tikzcd} \textbf{A} & B \end{tikzcd}
%%%% case macro-safe-redefined-mention-text
\renewcommand\itshape{Q}
\begin{tikzcd} \textit{A} \arrow[r] & \emph{B} \end{tikzcd}
%%%% case macro-safe-redefined-mention-row
\renewcommand\bfseries{\\}
\begin{tikzcd} \textbf{A} & B \end{tikzcd}
%%%% case macro-safe-redefined-mention-amp
\makeatletter\catcode`\&=13 \gdef\mdseries{&}\catcode`\&=4 \makeatother
\begin{tikzcd} \textmd{A} & B \end{tikzcd}
%%%% case macro-safe-mention-math
\renewcommand\mathbf{\arrow[r]}
\newcommand\h{\mathbf}
\begin{tikzcd} \h{A} & B \end{tikzcd}
%%%% case macro-safe-mention-amsmath-text
\renewcommand\ifmmode{}
\begin{tikzcd} A \arrow[r] & \text{B} \end{tikzcd}
%%%% case macro-safe-mention-protect
\let\protect\arrow
\begin{tikzcd} \textbf{A}[r] & B \end{tikzcd}
\let\protect\relax
%%%% case macro-safe-mention-operatorname
\makeatletter\renewcommand\qopname@{\arrow[r]}\makeatother
\begin{tikzcd} \operatorname{A} & B \end{tikzcd}
%%%% case macro-safe-mention-limits
\renewcommand\sum{\arrow[r]}
\begin{tikzcd} \sum & B \end{tikzcd}
%%%% case macro-safe-mention-hat
\renewcommand\mathaccent{\arrow[r]}
\begin{tikzcd} \hat{A} & B \end{tikzcd}
%%%% case macro-safe-mention-bar
\let\mathchoice\arrow
\begin{tikzcd} \phantom{A}[r] & B \end{tikzcd}
%%%% case macro-safe-primitive-let
\let\mathsf\relax
\begin{tikzcd} \mathsf{A} \arrow[r] & B \end{tikzcd}
%%%% case macro-safe-frac
\let\over\arrow
\begin{tikzcd} \frac{A}{B}[r] & B \end{tikzcd}
%%%% case macro-safe-relax
\let\oldrelax\relax
\let\relax\arrow
\begin{tikzcd} \mathbf{A}[r] & B \end{tikzcd}
%%%% case macro-safe-primitive-redefined
\let\hbox\arrow
\begin{tikzcd} \text{A}[r] & B \end{tikzcd}
%%%% case font-size-before
\large
\begin{tikzcd} A \arrow[r, "f"] & B \end{tikzcd}
%%%% case font-size-footnotesize
\footnotesize
\begin{tikzcd} A \arrow[r, "f"] & B \end{tikzcd}
%%%% case font-size-small
\small
\begin{tikzcd} A \arrow[r, "f"] & B \end{tikzcd}
%%%% case font-size-tiny
\tiny
\begin{tikzcd} A \arrow[r, "f"] & B \end{tikzcd}
%%%% case font-size-huge
\Huge
\begin{tikzcd} A \arrow[r, "f"] & B \end{tikzcd}
%%%% case font-size-units
\large
\begin{tikzcd}[row sep=1em, column sep=1ex] A \arrow[r, shift left=1ex, "f" {xshift=1em}] \arrow[d, shift left=1em] & B \\[1ex] C &[1em] D \end{tikzcd}
%%%% case font-size-unit-opt
\Large
\begin{tikzcd} A \arrow[d, "f"] \\[1em] B \end{tikzcd}
%%%% case font-small-opt-large
\small
\begin{tikzcd}[font=\large] A \arrow[r, "f"] & B \end{tikzcd}
%%%% case font-opt-units
\begin{tikzcd}[font=\large, row sep=1em, column sep=1ex] A \arrow[d, "f"] \\ B \end{tikzcd}
%%%% case font-opt-units-before
\begin{tikzcd}[row sep=1em, column sep=1ex, font=\large] A \arrow[d, "f"] \\ B \end{tikzcd}
%%%% case font-opt-units-in-cell
\begin{tikzcd}[font=\large] A \arrow[d, "f", shift left=1em] \\[1ex] B \end{tikzcd}
%%%% case font-opt-tiny
\begin{tikzcd}[font=\tiny] A \arrow[r, "f"] & B \end{tikzcd}
%%%% case font-opt-itshape
\begin{tikzcd}[font=\itshape] A \arrow[r, "f"] & B \end{tikzcd}
%%%% case font-opt-macro
\newcommand\q{\large}
\begin{tikzcd}[font=\q] A \arrow[r, "f"] & B \end{tikzcd}
%%%% case font-opt-macro-arrow
\newcommand\q{\arrow[r]}
\begin{tikzcd}[font=\q] A & B \end{tikzcd}
%%%% case font-opt-macro-text
\newcommand\q{Q}
\begin{tikzcd}[font=\q] A \arrow[r] & B \end{tikzcd}
%%%% case font-opt-value-text
\begin{tikzcd}[font=Q] A \arrow[r] & B \end{tikzcd}
%%%% case font-opt-empty
\begin{tikzcd}[font=] A \arrow[r] & B \end{tikzcd}
%%%% case font-opt-color
\begin{tikzcd}[font=\color{red}] A \arrow[r] & B \end{tikzcd}
%%%% case node-font-opt
\begin{tikzcd}[node font=\large] A \arrow[r, "f"] & B \end{tikzcd}
%%%% case every-matrix-font
\begin{tikzcd}[every matrix/.append style={font=\large}] A \arrow[r, "f"] & B \end{tikzcd}
%%%% case nodes-font
\begin{tikzcd}[nodes={font=\large}] A \arrow[r, "f"] & B \end{tikzcd}
%%%% case nodes-red
\begin{tikzcd}[nodes={red}] A \arrow[r, "f"] & B \end{tikzcd}
%%%% case nodes-draw
\begin{tikzcd}[nodes={draw}] A \arrow[r, "f"] & B \end{tikzcd}
%%%% case nodes-inner-sep
\begin{tikzcd}[nodes={inner sep=1pt}] A \arrow[r, "f"] & B \end{tikzcd}
%%%% case cells-empty-nodes
\begin{tikzcd}[cells={nodes={}}] A \arrow[r, "f"] & B \end{tikzcd}
%%%% case cells-text
\begin{tikzcd}[cells={nodes={text=red}}] A \arrow[r, "f"] & B \end{tikzcd}
%%%% case cells-font
\begin{tikzcd}[cells={nodes={font=\tiny}}] A \arrow[r, "f"] & B \end{tikzcd}
%%%% case cells-anchor
\begin{tikzcd}[cells={anchor=west}] A \arrow[r, "f"] & B \end{tikzcd}
%%%% case cells-text-height
\begin{tikzcd}[cells={nodes={text height=1cm}}] A \arrow[r, "f"] & B \end{tikzcd}
%%%% case cells-text-width
\begin{tikzcd}[cells={nodes={text width=1cm}}] A \arrow[r, "f"] & B \end{tikzcd}
%%%% case matrix-options-key
\begin{tikzcd}[matrix of nodes] A \arrow[r, "f"] & B \end{tikzcd}
%%%% case matrix-options-matrix
\begin{tikzcd}[execute at begin cell=Q] A \arrow[r, "f"] & B \end{tikzcd}
%%%% case matrix-options-end
\begin{tikzcd}[execute at end cell=Q] A \arrow[r, "f"] & B \end{tikzcd}
%%%% case matrix-options-empty
\begin{tikzcd}[execute at empty cell=Q] A \arrow[r, "f"] & \end{tikzcd}
%%%% case matrix-options-nodes-in-empty
\begin{tikzcd}[nodes in empty cells] A \arrow[r, "f"] & \end{tikzcd}
%%%% case matrix-options-nodes-in-empty-false
\begin{tikzcd}[nodes in empty cells=false] A \arrow[r, "f"] & \end{tikzcd}
%%%% case matrix-options-execute-picture
\begin{tikzcd}[execute at begin picture=\typeout{x}] A \arrow[r, "f"] & B \end{tikzcd}
%%%% case matrix-options-line-width
\begin{tikzcd}[line width=2pt] A \arrow[r, "f"] & |[draw]| B \end{tikzcd}
%%%% case matrix-options-scale
\begin{tikzcd}[scale=2] A \arrow[r, "f"] & B \end{tikzcd}
%%%% case matrix-options-transform-shape
\begin{tikzcd}[scale=2, transform shape] A \arrow[r, "f"] & B \end{tikzcd}
%%%% case matrix-options-rotate
\begin{tikzcd}[rotate=30] A \arrow[r, "f"] & B \end{tikzcd}
%%%% case matrix-options-xshift
\begin{tikzcd}[xshift=1cm] A \arrow[r, "f"] & B \end{tikzcd}
%%%% case matrix-options-inner-sep
\begin{tikzcd}[inner sep=1pt] A \arrow[r, "f"] & B \end{tikzcd}
%%%% case matrix-options-red
\begin{tikzcd}[red] A \arrow[r, "f"] & B \end{tikzcd}
%%%% case matrix-options-draw
\begin{tikzcd}[draw] A \arrow[r, "f"] & B \end{tikzcd}
%%%% case matrix-options-baseline
\begin{tikzcd}[baseline=1cm] A \arrow[r, "f"] & B \end{tikzcd}
%%%% case matrix-options-baseline-macro
\def\q{1cm}
\begin{tikzcd}[baseline=\q] A \arrow[r, "f"] & B \end{tikzcd}
%%%% case matrix-options-baseline-cell
\begin{tikzcd}[baseline=(\tikzcdmatrixname-1-1.base)] A \arrow[d, "f"] \\ B \end{tikzcd}
%%%% case matrix-options-baseline-center
\begin{tikzcd}[baseline=(current bounding box.center)] A \arrow[d, "f"] \\ B \end{tikzcd}
%%%% case matrix-options-baseline-alias
\begin{tikzcd}[baseline=(q.base)] |[alias=q]| A \arrow[d, "f"] \\ B \end{tikzcd}
%%%% case matrix-options-baseline-no-node
\begin{tikzcd}[baseline=(q.base)] A \arrow[d, "f"] \\ B \end{tikzcd}
%%%% case matrix-options-baseline-empty-cell
\begin{tikzcd}[baseline=(\tikzcdmatrixname-1-1.base)] \arrow[d, "f"] \\ B \end{tikzcd}
%%%% case matrix-options-remember
\begin{tikzcd}[remember picture] A \arrow[r, "f"] & B \end{tikzcd}
%%%% case matrix-options-name
\begin{tikzcd}[name=q] A \arrow[r, "f"] & B \end{tikzcd}
%%%% case matrix-options-column-sep-negative
\begin{tikzcd}[column sep=-1cm] A \arrow[r, "f"] & B \end{tikzcd}
%%%% case matrix-options-row-sep-origins
\begin{tikzcd}[row sep={2cm,between origins}] A \arrow[d, "f"] \\ B \end{tikzcd}
%%%% case matrix-options-column-sep-borders
\begin{tikzcd}[column sep={2cm,between borders}] A \arrow[r, "f"] & B \end{tikzcd}
%%%% case matrix-options-seps-macro
\def\q{2cm,between origins}
\begin{tikzcd}[column sep=\q] A \arrow[r, "f"] & B \end{tikzcd}
%%%% case matrix-options-sep-expression
\begin{tikzcd}[column sep=1cm+2pt, row sep=2*3pt] A \arrow[r, "f"] & B \end{tikzcd}
%%%% case matrix-options-sep-function
\begin{tikzcd}[column sep=sqrt(2)*1cm] A \arrow[r, "f"] & B \end{tikzcd}
%%%% case matrix-options-sep-dimexpr
\begin{tikzcd}[column sep=\dimexpr 1cm+2pt\relax] A \arrow[r, "f"] & B \end{tikzcd}
%%%% case matrix-options-sep-length
\newlength\q \setlength\q{1cm}
\begin{tikzcd}[column sep=\q] A \arrow[r, "f"] & B \end{tikzcd}
%%%% case matrix-options-sep-length-times
\newlength\q \setlength\q{1cm}
\begin{tikzcd}[column sep=2\q] A \arrow[r, "f"] & B \end{tikzcd}
%%%% case matrix-options-sep-count
\newcount\q \q=20
\begin{tikzcd}[column sep=\q pt] A \arrow[r, "f"] & B \end{tikzcd}
%%%% case matrix-options-sep-the
\newlength\q \setlength\q{1cm}
\begin{tikzcd}[column sep=\the\q] A \arrow[r, "f"] & B \end{tikzcd}
%%%% case matrix-options-sep-baselineskip
\begin{tikzcd}[row sep=\baselineskip] A \arrow[d, "f"] \\ B \end{tikzcd}
%%%% case matrix-options-sep-hsize
\begin{tikzcd}[column sep=.1\textwidth] A \arrow[r, "f"] & B \end{tikzcd}
%%%% case matrix-options-sep-linewidth
\begin{tikzcd}[column sep=.1\linewidth] A \arrow[r, "f"] & B \end{tikzcd}
%%%% case matrix-options-sep-unit-ex
\begin{tikzcd}[column sep=2ex, row sep=1.5em] A \arrow[r, "f"] \arrow[d] & B \\ C \end{tikzcd}
%%%% case matrix-options-sep-unit-other
\begin{tikzcd}[column sep=2mm, row sep=1bp] A \arrow[r, "f"] \arrow[d] & B \\ C \end{tikzcd}
%%%% case matrix-options-sep-unit-more
\begin{tikzcd}[column sep=2dd, row sep=1cc] A \arrow[r, "f"] \arrow[d] & B \\ C \end{tikzcd}
%%%% case matrix-options-sep-unit-sp-in-nd
\begin{tikzcd}[column sep=2000000sp, row sep=1in] A \arrow[r, "f"] \arrow[d] & B \\ C \end{tikzcd}
%%%% case matrix-options-sep-unit-pc
\begin{tikzcd}[column sep=2pc, row sep=1pc] A \arrow[r, "f"] \arrow[d] & B \\ C \end{tikzcd}
%%%% case matrix-options-sep-no-unit
\begin{tikzcd}[column sep=20, row sep=10] A \arrow[r, "f"] \arrow[d] & B \\ C \end{tikzcd}
%%%% case matrix-options-sep-true
\begin{tikzcd}[column sep=2 truecm] A \arrow[r, "f"] & B \end{tikzcd}
%%%% case matrix-options-sep-true-mag
\mag=1200
\begin{tikzcd}[column sep=2 truecm] A \arrow[r, "f"] & B \end{tikzcd}
%%%% case matrix-options-sep-int
\begin{tikzcd}[column sep=2 * 1cm] A \arrow[r, "f"] & B \end{tikzcd}
%%%% case font-parbox-in-cell
\begin{tikzcd} {\large A} \arrow[r] & {\fontsize{14}{16}\selectfont B} \end{tikzcd}
%%%% case font-fontsize-in-cell
\begin{tikzcd} \fontsize{14}{16}\selectfont A \arrow[r] & \large B \end{tikzcd}
%%%% case font-fontsize-before
\fontsize{13}{15}\selectfont
\begin{tikzcd} A \arrow[r, "f"] & B \end{tikzcd}
%%%% case font-fontsize-small-unit
\fontsize{7.3}{9}\selectfont
\begin{tikzcd}[column sep=1em] A \arrow[r, "f"] & B \end{tikzcd}
%%%% case font-mathversion-bold
\mathversion{bold}
\begin{tikzcd} A \arrow[r, "f"] & B \end{tikzcd}
%%%% case font-mathversion-bold-in-cell
\begin{tikzcd} \mathversion{bold} A \arrow[r, "f"] & B \end{tikzcd}
%%%% case font-boldmath
\boldmath
\begin{tikzcd} A \arrow[r, "f"] & B \end{tikzcd}
%%%% case font-boldmath-in-cell
\begin{tikzcd} \boldmath A \arrow[r, "f"] & \unboldmath B \end{tikzcd}
%%%% case font-bfseries-before
\bfseries
\begin{tikzcd} A \arrow[r, "f"] & \text{word} \end{tikzcd}
%%%% case font-itshape-before
\itshape
\begin{tikzcd}[column sep=1em] A \arrow[r, "f"] & \text{word} \end{tikzcd}
%%%% case font-sffamily-before
\sffamily
\begin{tikzcd}[column sep=1em] A \arrow[r, "f"] & \text{word} \end{tikzcd}
%%%% case font-ttfamily-before
\ttfamily
\begin{tikzcd}[column sep=1em] A \arrow[r, "f"] & \text{word} \end{tikzcd}
%%%% case font-font-primitive-before
\font\q=cmr17 \q
\begin{tikzcd}[column sep=1em] A \arrow[r, "f"] & \text{word} \end{tikzcd}
%%%% case font-font-primitive-in-cell
\font\q=cmr17
\begin{tikzcd} \text{\q word} \arrow[r] & \q B \end{tikzcd}
%%%% case font-nullfont-before
\nullfont
\begin{tikzcd}[column sep=1em] A \arrow[r, "f"] & \text{word} \end{tikzcd}
%%%% case font-mathfont-changed
\textfont0=\font
\begin{tikzcd} \mathrm{A} \arrow[r, "f"] & 1 \end{tikzcd}
%%%% case font-rule-thickness-changed
\fontdimen8\textfont3=2pt
\begin{tikzcd} A \arrow[r, "f"] & B \end{tikzcd}
%%%% case font-axis-height-changed
\fontdimen22\textfont2=5pt
\begin{tikzcd} A \arrow[r, "f"] & B \end{tikzcd}
%%%% case font-fontdimen-text
\fontdimen5\font=10pt \fontdimen6\font=30pt
\begin{tikzcd}[column sep=1em, row sep=1ex] A \arrow[r, "f"] \arrow[d] & B \\ C \end{tikzcd}
%%%% case font-setmathalphabet
\SetMathAlphabet{\mathrm}{normal}{OT1}{cmss}{m}{n}
\begin{tikzcd} \mathrm{A} \arrow[r] & 1 \end{tikzcd}
%%%% case font-setsymbolfont
\SetSymbolFont{operators}{normal}{OT1}{cmss}{m}{n}
\begin{tikzcd} \mathrm{A} \arrow[r] & 1 \end{tikzcd}
%%%% case font-sizes-math
\DeclareMathSizes{10}{20}{15}{10}
\begin{tikzcd} A^b \arrow[r, "f"] & 1 \end{tikzcd}
%%%% case lmodern
%%%% pre
\usepackage{lmodern}
%%%% body
\begin{tikzcd}[column sep=1em] A \arrow[r, "f"] & \text{word} \end{tikzcd}
%%%% case lmodern-t1
%%%% pre
\usepackage[T1]{fontenc}\usepackage{lmodern}
%%%% body
\begin{tikzcd}[column sep=1em] A \arrow[r, "f", hook] & \text{word} \end{tikzcd}
%%%% case times
%%%% pre
\usepackage{times}
%%%% body
\begin{tikzcd}[column sep=1em] A \arrow[r, "f"] & \text{word} \end{tikzcd}
%%%% case mathptmx
%%%% pre
\usepackage{mathptmx}
%%%% body
\begin{tikzcd}[column sep=1em] A \arrow[r, "f"] \arrow[d, Rightarrow] & \text{word} \\ C \end{tikzcd}
%%%% case mathpazo
%%%% pre
\usepackage{mathpazo}
%%%% body
\begin{tikzcd}[column sep=1em] A \arrow[r, "f"] \arrow[d, Rightarrow] & \text{word} \\ C \end{tikzcd}
%%%% case newtx
%%%% pre
\usepackage{newtxtext,newtxmath}
%%%% body
\begin{tikzcd}[column sep=1em] A \arrow[r, "f"] \arrow[d, Rightarrow] & \text{word} \\ C \end{tikzcd}
%%%% case newpx
%%%% pre
\usepackage{newpxtext,newpxmath}
%%%% body
\begin{tikzcd}[column sep=1em] A \arrow[r, "f"] \arrow[d, Rightarrow] & \text{word} \\ C \end{tikzcd}
%%%% case stix
%%%% pre
\usepackage{stix}
%%%% body
\begin{tikzcd}[column sep=1em] A \arrow[r, "f"] \arrow[d, Rightarrow] & \text{word} \\ C \end{tikzcd}
%%%% case fourier
%%%% pre
\usepackage{fourier}
%%%% body
\begin{tikzcd}[column sep=1em] A \arrow[r, "f"] \arrow[d, Rightarrow] & \text{word} \\ C \end{tikzcd}
%%%% case eulervm
%%%% pre
\usepackage{eulervm}
%%%% body
\begin{tikzcd}[column sep=1em] A \arrow[r, "f"] \arrow[d, Rightarrow] & \text{word} \\ C \end{tikzcd}
%%%% case cmbright
%%%% pre
\usepackage{cmbright}
%%%% body
\begin{tikzcd}[column sep=1em] A \arrow[r, "f"] \arrow[d, Rightarrow] & \text{word} \\ C \end{tikzcd}
%%%% case microtype
%%%% pre
\usepackage{microtype}
%%%% body
\begin{tikzcd}[column sep=1em] A \arrow[r, "f"] \arrow[d, Rightarrow] & \text{word} \\ C \end{tikzcd}
Some text that is long enough to fill up a whole line in this document and break.
%%%% case geometry
%%%% pre
\usepackage[margin=1cm]{geometry}
%%%% body
\begin{tikzcd}[column sep=1em] A \arrow[r, "f"] & \text{word} \end{tikzcd}
%%%% case xcolor-svgnames
%%%% pre
\usepackage[svgnames]{xcolor}
%%%% body
\begin{tikzcd} A \arrow[r, Crimson] & B \end{tikzcd}
%%%% case xcolor-cmyk
%%%% pre
\usepackage[cmyk]{xcolor}
%%%% body
\begin{tikzcd} A \arrow[r, red] & |[draw=blue]| B \end{tikzcd}
%%%% case xcolor-gray
%%%% pre
\usepackage[gray]{xcolor}
%%%% body
\begin{tikzcd} A \arrow[r, red] & B \end{tikzcd}
%%%% case xcolor-rgb-target
%%%% pre
\usepackage[rgb]{xcolor}
%%%% body
\begin{tikzcd} A \arrow[r, cyan] & \textcolor{cyan}{B} \end{tikzcd}
%%%% case xcolor-natural
%%%% pre
\usepackage[natural]{xcolor}
%%%% body
\begin{tikzcd} A \arrow[r, cyan] & \textcolor{cyan}{B} \end{tikzcd}
%%%% case amscd-too
%%%% pre
\usepackage{amscd}
%%%% body
\begin{tikzcd} A \arrow[r] & B \end{tikzcd}
\begin{CD} A @>>> B \end{CD}
%%%% case xypic-too
%%%% pre
\usepackage[all]{xy}
%%%% body
\begin{tikzcd} A \arrow[r] & B \end{tikzcd}
\xymatrix{A \ar[r] & B}
%%%% case breqn-too
%%%% pre
\usepackage{breqn}
%%%% body
\begin{tikzcd} A \arrow[r] & B \end{tikzcd}
%%%% case tikz-after
%%%% pre
\usepackage{tikz}\usetikzlibrary{arrows.meta,calc,decorations.pathreplacing,positioning}
%%%% body
\begin{tikzcd} A \arrow[r] & B \end{tikzcd}
%%%% case tikz-nfold
%%%% pre
\usetikzlibrary{nfold}
%%%% body
\begin{tikzcd} A \arrow[r, Rightarrow] \arrow[d, equal] & B \\ C \end{tikzcd}
%%%% case tikz-external
%%%% pre
\usetikzlibrary{external}
\tikzexternalize
%%%% body
\begin{tikzcd} A \arrow[r] & B \end{tikzcd}
%%%% case tikz-backgrounds-layers
%%%% pre
\usetikzlibrary{backgrounds}
%%%% body
\begin{tikzcd} A \arrow[r] & B \end{tikzcd}
\begin{tikzpicture}\begin{scope}[on background layer]\node {x};\end{scope}\end{tikzpicture}
%%%% case tikz-layers-set
\pgfdeclarelayer{back}
\pgfsetlayers{back,main}
\begin{tikzcd} A \arrow[r] & B \end{tikzcd}
%%%% case tikz-layers-set-no-main
\pgfdeclarelayer{back}
\pgfsetlayers{back}
\begin{tikzcd} A \arrow[r] & B \end{tikzcd}
%%%% case quiver
%%%% pre
\usepackage{quiver}
%%%% body
\begin{tikzcd} A \arrow[r, "\shortmid"{marking}, curve={height=6pt}] & B \end{tikzcd}
